# Title
Fish: A Review of Classification, Biodiversity, and Research Gaps

## Abstract
This paper provides an overview of the **definition and classification of fish**, highlighting their diversity and the challenges associated with estimating the number of species. The analysis reveals discrepancies in species counts, with estimates ranging from 20,000 to 34,000, and identifies key gaps in research, including **taxonomic uncertainty**, **lack of academic sources**, and **geographic and ecological gaps**. The study recommends the use of standardized taxonomic databases, incorporation of peer-reviewed studies, and development of a unified framework for classifying and tracking fish biodiversity.

## 1. Introduction
Fish are a diverse group of **aquatic, anamniotic, gill-bearing vertebrate animals** with a cranium but no limbs with digits [1]. They are classified into approximately 20,000–34,000 species, according to sources such as **NOAA Fisheries** and **Britannica**. This paper aims to review the current state of knowledge on fish, including their classification, biodiversity, and notable facts, as well as identify research gaps and provide recommendations for future studies.

## 2. Related Work
Previous studies have highlighted the challenges associated with estimating the number of fish species, with discrepancies in species counts reflecting differences in methodological scope [2]. **General sources**, such as **Wikipedia** and **Britannica**, provide broad, encyclopedic definitions and classifications, while **specialized sources**, such as **NOAA Fisheries** and **Predatory Fins**, offer targeted facts and tools for enthusiasts. However, the lack of **academic sources** limits deeper insights into evolutionary biology, ecology, or conservation.

## 3. Methodology
This study is based on a review of existing literature and online sources, including **Wikipedia**, **Britannica**, **NOAA Fisheries**, **Predatory Fins**, and **FishMap.org**. The analysis focuses on the **definition and classification of fish**, **biodiversity and species estimates**, and **notable facts and extremes**. The study also identifies **research gaps** and provides **recommendations** for future studies.

## 4. Results & Findings
The analysis reveals that:
* Fish are a diverse group of approximately 20,000–34,000 species
* Discrepancies exist in species counts, with estimates ranging from 20,000 to 34,000
* The **whale shark** is the largest fish species
* **Platforms like Predatory Fins and FishMap.org** provide specialized information on freshwater species and geographic distributions
* Key gaps in research include **taxonomic uncertainty**, **lack of academic sources**, and **geographic and ecological gaps**

## 5. Discussion
The findings of this study highlight the need for standardized taxonomic databases and the incorporation of peer-reviewed studies to explore ecological, evolutionary, or conservation-focused questions. The development of a unified framework for classifying and tracking fish biodiversity is also recommended. Furthermore, the study suggests that future research should narrow its scope to address open problems meaningfully, such as focusing on freshwater or marine fish, or specific families.

## 6. Conclusion
In conclusion, this paper provides an overview of the current state of knowledge on fish, including their classification, biodiversity, and notable facts. The study identifies key gaps in research and provides recommendations for future studies, including the use of standardized taxonomic databases, incorporation of peer-reviewed studies, and development of a unified framework for classifying and tracking fish biodiversity. By addressing these gaps, future research can contribute to a better understanding of fish and their role in aquatic ecosystems.

## References
[1] Wikipedia, Definition of Fish, 2022
[2] Britannica, Fish, 2022
[3] NOAA Fisheries, Fish Species, 2022
[4] Predatory Fins, Freshwater Fish, 2022
[5] FishMap.org, Geographic Distribution of Fish, 2022